\documentclass[12pt,a4paper]{ltjsarticle}
\usepackage[margin=22mm]{geometry}
\usepackage{xcolor}
\usepackage{luatexja-fontspec}
\setlength{\parindent}{0pt}
\setlength{\parskip}{0.7em}
\setmainjfont{HaranoAjiMincho-Regular}
\setsansjfont{HaranoAjiGothic-Medium}
\linespread{1.18}\selectfont
\ltjsetparameter{kanjiskip=0.20em plus 0.10em minus 0.06em}
\ltjsetparameter{xkanjiskip=0.24em plus 0.10em minus 0.06em}

\newcommand{\sw}[1]{\textbf{#1}}

\begin{document}

\begin{center}
{\LARGE Presentation1 Script}
\end{center}

\sw{H\.{e}llo} everyone. $\downarrow$ //

My name is \sw{H\.{e}izo N\.{a}kai}, / and my presenter number is \sw{55}. $\downarrow$ //

\sw{F\.{i}rst}, / I'd like to introduce \sw{mys\.{e}lf}. $\downarrow$ //

I'm from \sw{Sh\.{i}ga Pr\.{e}fecture}, / and I lived there until I graduated from \sw{h\.{i}gh school}. $\downarrow$ //

\sw{Sh\.{i}ga} is famous for \sw{L\.{a}ke B\.{i}wa}, / the largest lake in \sw{Jap\.{a}n}. $\downarrow$ //

I grew up in a \sw{c\.{a}lm env\.{i}ronment} / surrounded by \sw{n\.{a}ture}. $\downarrow$ //

After \sw{h\.{i}gh school}, / I moved to \sw{T\.{o}kyo} for university. $\downarrow$ //

During my undergraduate years, / I became interested in \sw{t\.{e}chnology}, / especially \sw{v\.{i}rtual r\.{e}ality} and \sw{h\.{u}man-comp\.{u}ter interact\.{i}on}. $\downarrow$ //

\sw{N\.{o}w}, / I'm a first-year \sw{m\.{a}ster's st\.{u}dent} here at \sw{N\.{A}IST}. $\downarrow$ //

\sw{N\.{e}xt}, / I'd like to talk about my \sw{res\.{e}arch field}. $\downarrow$ //

I belong to the \sw{Comp\.{u}ting Arch\.{i}tecture Laboratory}, / and my research field is \sw{comp\.{u}ting arch\.{i}tecture}. $\downarrow$ //

At first, / I was mainly interested in \sw{v\.{i}rtual r\.{e}ality}. $\downarrow$ //

\sw{How\.{e}ver}, / as I learned more about \sw{V\.{R} systems}, $\searrow$ / I realized that \sw{high-qu\.{a}lity V\.{R}} requires a lot of \sw{comput\.{a}tion} and \sw{\.{e}nergy}. $\downarrow$ //

For example, / VR needs \sw{real-time \.{i}mage processing}, / \sw{low l\.{a}tency}, / and \sw{high perf\.{o}rmance}. $\downarrow$ //

This made me interested in the \sw{h\.{a}rdware} and \sw{s\.{y}stem-level techn\.{o}logies} / behind these applications. $\downarrow$ //

I think \sw{comp\.{u}ting arch\.{i}tecture} is interesting / because it connects \sw{h\.{a}rdware} and \sw{s\.{o}ftware}. $\downarrow$ //

By improving \sw{comp\.{u}ter arch\.{i}tecture}, $\searrow$ / we can make advanced applications \sw{f\.{a}ster} / and more \sw{\.{e}nergy eff\.{i}cient}. $\downarrow$ //

In the \sw{f\.{u}ture}, / I would like to become an \sw{engin\.{e}er} or \sw{res\.{e}archer} / who understands both \sw{h\.{a}rdware} and \sw{s\.{o}ftware}. $\downarrow$ //

Thank you for \sw{l\.{i}stening}. $\downarrow$ //

\section*{Delivery II: Speed, pronunciation, \& stress p.2（日本語訳）}

\textbf{1. どうすれば、話すスピードが速くなりすぎるのを防げるか？}

\begin{itemize}
  \item 緊張をコントロールする：リハーサルとイメージトレーニングを行う
  \item 専門用語や重要なフレーズは、よりゆっくり・はっきり話す
\end{itemize}

\textbf{2. どのような場面で間を取るべきか？}

\begin{itemize}
  \item 短い間：節と節、文と文のあいだの文法的な切れ目
  \item 長い間：重要なポイントの後、修辞疑問文の後、話の段階の切り替え時、スライドを切り替える前
\end{itemize}

\textbf{3. 発表の中で、どの語の発音に特に注意すべきか？}

\begin{itemize}
  \item キーワード：繰り返し出てくる語、主要ポイントに関連する専門語
\end{itemize}

\end{document}
